\labsection{4}{Circular linked-list operations}{sec:lab04}

\begin{sourcealignment}{Official Lab Manual --- Lab IV}
\textbf{Objective.} Study linked-list types more deeply and implement algorithms for circular linked lists.

\textbf{Outcomes.} Understand circular linked lists and implement one from scratch.

\textbf{Problem directions.} Implement the different insertion/deletion operations possible in a circular linked list; compare singly, doubly, and circular lists.
\end{sourcealignment}

\begin{labproject}{Maintain a circular list without breaking the loop}
\textbf{Coverage.} Chapter 4.\\
\textbf{Goal.} Maintain a circular list through every insertion and deletion without breaking the ring.

Implement insertion at beginning/end, deletion from beginning/end, deletion by value, and display.

\begin{lstlisting}
void displayCircular(Node* last) {
    if (last == nullptr) return;
    Node* head = last->next;
    Node* p = head;
    do {
        cout << p->data << ' ';
        p = p->next;
    } while (p != head);
}
\end{lstlisting}

\begin{tryit}
Use the same input sequence in singly, doubly, and circular implementations. Compare traversal stop conditions and the pointer fields required per node.
\end{tryit}
\end{labproject}

\begin{verification}
\begin{itemize}
  \item Empty, one-node, and multi-node cases are tested.
  \item After every mutation, the last node still points to the head.
  \item Traversal stops when it returns to the starting node rather than waiting for \texttt{nullptr}.
  \item Deletion-by-value handles the head, a middle node, and the last node.
\end{itemize}
\end{verification}

\begin{labdeliverable}
Submit the required operations, test output, and a three-column comparison of singly, doubly, and circular linked lists.
\end{labdeliverable}
